\documentclass[../main.tex]{subfiles}

\begin{document}
  \begin{tikzpicture}[remember picture, overlay]
    % Education
    \node (education) at ($(current page.north west)+(\paperwidth-\cvMiddleSidebar+\Margin / 2,\textStart)+(-0.5mm,-1.8\Margin+1.8cm)$)
      [rectangle, fill=education, minimum width=\paperwidth-\cvMiddleSidebar, minimum height=\cvSectionHeight]
      {\Large \textbf{\textcolor{text}{\faGraduationCap\space Edukacja}}};
    \node (education data) at ($(education.south) + (0cm, -2.2mm)$) [below, text width = \paperwidth-\cvMiddleSidebar-\Margin] {
      \educationFinished{\href{https://mba.agh.edu.pl/programs/mba-executive}{Executive MBA}}
      {\href{https://www.agh.edu.pl/en/}{AGH}}

      \vspace*{3.7mm}
      \educationFinished{Informatyka, studia magisterskie}
      {\href{https://en.uj.edu.pl/}{Uniwersytet Jagielloński}}
    };

    % Work
    \node (work) at ($(education data.south)+(0cm, -2.35mm)$)
      [rectangle, fill=work, below,
      minimum width=\paperwidth-\cvMiddleSidebar, minimum height=\cvSectionHeight]
      {\Large \textbf{\textcolor{text}{\faBriefcase\space Praca}}};
    \node (work data) at (work.south) [below, text width = \paperwidth-\cvMiddleSidebar-\Margin] {
      \begin{cvtext}
        \work{Kandydat na CTO}{\href{https://softwise.ai/}{softwise.ai}}{10/2025}{Obecnie}
        {
          {
            {ulepszenie wieloagentowego systemu (5× niższym opóźnieniem)},
            {zarządzanie zespołem i sprzedażą od strony technicznej (pre-sales)}
          }
        }
        {
          {analiza wykonawców, SI, LLM, programowanie, zarządzanie, sprzedaż}
        }
        \work{Założyciel}{\href{https://www.opennudge.com}{opennudge}}{04/2024}{Obecnie}
        {
          {
            {prowadzenie organizacji tworzącej oprogramowanie open-source},
            {tworzenie narzędzi upraszczających tworzenie oprogramowania}
          }
        }
        {
          {FOSS, zgodność, SI, Dev(Sec)Ops, programowanie, agenci, zarządzanie}
        }
        \work{CTO}{\href{https://www.inovintell.com}{inovintell}}{04/2022}{02/2024}
        {
          {
            założenie firmy AI–life science wraz z interesariuszami wewnętrznymi,
            kierowanie i mentoring zespołów technicznych i badawczych (10 osób),
            {tworzenie strategii projektowej, technologicznej i kultury firmowej}
          }
        }
        {
          {life science, SI, ML, DS, DL, NLP, LLM, programowanie, zarządzanie}
        }
        \work{Doradca sztucznej inteligencji}{\href{https://www.wincan.com/en/home/}{wincan}}{10/2022}{06/2023}
        {
          {
            mentoring oraz doradztwo strategiczne dla zespołu machine learning,
            wsparcie rozwoju algorytmów SI (wizja komputerowa)
          }
        }{}
        \work{Doradca sztucznej inteligencji}{\href{https://music.ai/}{music.ai}}{05/2022}{07/2022}
        {
          {
            optymalizacja sieci neuronowych pod wdrożenia na urządzenia mobilne,
            doradztwo w zakresie konwersji modeli sztucznej inteligencji
          }
        }{}
        \work{Szef programu nauczania}{\href{https://www.theaicore.com/}{AiCore}}{11/2020}{09/2021}
        {
          {
            projektowanie i prowadzenie modułów kursów Python{,} Linux i AI,
            wielokrotnie wybierany przez studentów jako najlepszy instruktor firmy,
            nadzór nad rozwojem technicznym i projektami ponad 100 studentów
          }
        }{}
        \work{Programista machine learning}{\href{https://pl.pinterest.com/artplanet_store/}{ArtPlanet}}{08/2019}{06/2020}
        {
          {
            zaprojektowanie niskokosztowej sieci kategoryzującej sztukę,
            odpowiedzialność za cały SDLC w tym DS oraz ML{,} wdrożenie na AWS
          }
        }{}
        \work{Programista machine learning}{\href{https://codete.com/}{Codete}}{04/2018}{10/2018}
        {
          {
            budowa i walidacja prototypu konwertera sieci neuronowych,
            współautor komercyjnych kursów machine learning firmy,
          }
        }{}
      \end{cvtext}
      \vspace*{-0.6cm}
      \plclause
    };

  \end{tikzpicture}
\end{document}
